\section{POTWORY, NEGOCJACJE I WYKUP}

Ogólna zasada:

W trakcie gry oddział często będzie stawać twarzą w twarz z potwornymi mieszkańcami labiryntu.
W takim przypadku będzie miał do wyboru trzy rodzaje zachowań: negocjacje, próbę wykupienia się lub walkę.
Zakończenie sukcesem negocjacje pozwala oddziałowi kontynuować eksplorację labiryntu bez konieczności toczenia walki z potworem, na którego się akurat natknął.
Jeżeli negocjacje się nie powiodą oddział może jeszcze próbować się wykupić.
Jeżeli i to się nie uda pozostaje tylko walka.
Jej wynikiem może być jedynie śmierć którejś z walczących stron.
Jeżeli oddział zwycięży, należące do niego postacie będą mogły podzielić między siebie posiadane przez pokonanego potwora skarby.

Sposób postępowania:

Gdy oddział wchodzi do nowego segmentu musi sprawdzić, czy w tym segmencie są potwory.
Sprawdzanie odbywa się przez rzut 1K6.

Jeżeli nowy segment jest komnatą, której oddział jeszcze nie odwiedzał, wyrzucenie 1,2 lub 3 oczek oznacza, że w tej komnacie ma swoje legowisko potwór osiadły.
W tabeli 8.3 ,,Potwory osiadłe'' należy sprawdzić (rzucając dwa razy 1K6) jego rodzaj.

Jeżeli nowym segmentem jest komnata, w której oddział już był lub jakikolwiek korytarz wyrzucenie 1 oczka oznacza, iż oddział natknął się na potwora wędrownego.
Należy sprawdzić jego rodzaj w tabeli 8.3 ,,Potwory wędrowne''.

Zasady szczegółowe:

[8.1] Jeżeli rzut kostką wskazuje, że nastąpiło spotkanie z potworem należy sprawdzić jego rodzaj w odpowiedniej do typu potwora (osiadły lub wędrowny) tabeli.

Jeden z graczy rzuca dwa razy 1K6 i na przecięciu rzędu, odpowiadającego liczbie oczek, wyrzuconych za pierwszym razem oraz kolumny, odpowiadającej liczbie oczek, wyrzuconych przy drugim rzucie odczytuje rodzaj oraz liczbę napotkanych potworów.
Należy przy tym pamiętać o tym, by uwzględnić poziom labiryntu~-- patrz 6.8.

[8.2] tabela ,,Charakterystyka potworów''~-- patrz plansza z tabelami.

[8.3] tabele: ,,Potwory osiadłe'', ,,Potwory wędrowne''~-- patrz plansza z tabelami.

[8.4] Gracze mogą podjąć próbę negocjacji z każdym potworem z wyjątkiem Bezimiennego i Demonów.

Mechanizm negocjacji jest następujący:

a. W tabeli 8.2 odnajdowany jest współczynnik ,,niechęć do negocjacji'' napotkanego potwora lub grupy potworów tego samego gatunku.

b. Jeden z graczy rzuca 2K6 i od sumy wyrzuconych oczek odejmuje ustalony w punkcie a. współczynnik, zaś do otrzymanego wyniku dodaje wielkość współczynnika zdolności ,,negocjacje'' dowolnego członka oddziału, np. zdolności ,,negocjacje'' równe +2 oznaczają, że do sumy oczek, wyrzuconych 2K6 należy dodać 2.
Na wynik negocjacji mogą mieć również wpływ odpowiednie czary (patrz 10.7).

c. Otrzymany w punkcie b. ostateczny wynik odnajdowany jest w tabeli 8.5 ,,Negocjacje'' i odczytywany jest odpowiadający mu rezultat negocjacji.

[8.5] Tabela ,,Negocjacje''~-- patrz plansza z tabelami.

[8.6] W tabeli 8.5 umieszczone są trzy możliwe rezultaty negocjacji:

niepowodzenie~-- podjęta próba nie powiodła się i oddział może usiłować się wykupić lub też podjąć z potworem walkę,

porozumienie~-- jeżeli oddział nie zaatakuje potwora, potwór również nie zaatakuje oddziału.
Oddział może przystąpić do oględzin znalezisk, jeżeli takie znajdują się w danym segmencie lub też pójść dalej.

zastraszenie~-- oddziałowi udaje się przekonać potwora, że nie ma on najmniejszych szans w bezpośrednim starciu, w związku z czym potwór nie tylko zostawia oddział w spokoju, ale również oddaje mu jedną czwartą posiadanych skarbów jako wykup.
Po wzięciu tej jednej czwartej oddział może oczywiście zaatakować potwora (np. by zdobyć cały skarb).

[8.7] Jeżeli negocjacje zawiodą oddział może usiłować kupić wolność, ofiarowując potworowi (nie dotyczy to Bezimiennego i Demonów) określoną sumę sztuk złota.
Sprawdzenie czy wykup został przyjęty przebiega następująco:

a. oblicza się sumę współczynników siła i niechęć do negocjacji najsilniejszego potwora w grupie,

b. przeznacza się określoną sumę sztuk złota jako wykup,

c. w tabeli 8.9 ,,Wykup'' na przecięciu kolumny, odpowiadającej obliczonej w punkcie a. sumie oraz rzędu, odpowiadającego ofiarowanej liczbie sztuk złota odnajdowana jest cyfra od 0 do 6,

d. jeden z graczy rzuca 1K6 i porównuje wyrzuconą liczbę oczek z odczytaną w punkcie c. cyfrą.
Jeżeli wyrzucona liczba oczek jest mniejsza lub równa tej cyfrze wykup został przyjęty.

Przyjęcie przez potwora wykupu oznacza dla oddziału sytuację analogiczną do tej, w której znalazłby się on, gdyby w negocjacjach osiągnięty został wynik ,,porozumienie'' (patrz 8.6).
Jeżeli potwór nie przyjmie wykupu atakuje on oddział jako pierwszy (opuszczana jest faza ataku oddziału).
Bez względu na to, czy potwór zaakceptuje wykup, czy też nie oddział traci zaofiarowane sztuki złota.
Oczywiście jeżeli potwór zostanie zabity pieniądze te mogą być mu odebrane z powrotem, razem z jego własnym skarbem.

[8.8] Wobec każdego potwora lub grupy potworów, bez względu na jej liczebność, może być podjęta tylko jedna próba wykupu.

[8.9] tabela ,,Wykup''~-- patrz plansza z tabelami.
